\section{Methodology}
\label{sec:method}

Given an RGB image with known camera intrinsics $\K$ and a set of known objects, our goal is to
predict, for every visible object instance, its rotation $\R \in SO(3)$ and translation
$\T \in \mathbb{R}^3$ from the model frame to the camera frame. As in
FastPose-ViT~\cite{ancey2026fastposevit}, a 2D detector first provides a bounding box and an
object label for each instance; the pose network then works on the image crop around each box.
\Cref{sec:method:recap} summarizes the original formulation, and
\cref{sec:method:translation,sec:method:conditioning,sec:method:symmetry,sec:method:training}
describe what changes for multiple objects and cameras.

\subsection{Recap: FastPose-ViT}
\label{sec:method:recap}
The crop is resized to the input size of a ViT-B/16~\cite{dosovitskiy2021vit} pretrained with
SWAG~\cite{singh2022swag}, and a single linear layer on the [CLS] token outputs nine values: three
normalized translation parameters $\U = (U_x, U_y, U_z)$ and a 6D rotation
representation~\cite{zhou2019continuity}, converted into a rotation matrix by Gram--Schmidt
orthonormalization. The network does not regress the absolute rotation $\R$, which depends on
where the object lies in the image, but the \emph{apparent rotation}
$\R' = \Delta\R\,\R$, where $\Delta\R$ rotates the viewing ray of the object centre
$\hat{\T} = \T / \|\T\|$ onto the optical axis $\mathbf{e}_z$. $\Delta\R$ has axis
$\mathbf{u} = \hat{\T} \times \mathbf{e}_z$ and angle $\theta = \arccos(\hat{\T}^\top
\mathbf{e}_z)$, and is computed in closed form with Rodrigues' formula; at inference,
$\R = \Delta\R^\top \R'$ once $\T$ is known. We keep the apparent rotation, the 6D
representation and the relative Frobenius rotation loss of~\cite{ancey2026fastposevit} unchanged,
and refer to it for their derivation.

\subsection{Camera- and size-independent translation}
\label{sec:method:translation}
Let $(b_x, b_y)$, $w$ and $h$ be the centre, width and height of the bounding box in pixels, and
$f_x, f_y, c_x, c_y$ the intrinsics of the image, which on BOP change from one image (and dataset)
to the next. The lateral targets measure where the projected object centre lies relative to the
box, in box units, as the scale-invariant translation of CDPN~\cite{li2019cdpn}:
\begin{equation}
  U_x = \frac{f_x \, t_x / t_z + c_x - b_x}{w}, \qquad
  U_y = \frac{f_y \, t_y / t_z + c_y - b_y}{h}.
  \label{eq:uxy}
\end{equation}
For the depth, the original formulation used a crop-depth proxy proportional to $t_z$ times the
box size divided by the image size. This proxy is well scaled for a single object and a single
camera, but on BOP it depends on the object size and the focal length: on LM-O it ranges from
about 50 to 300 (with $t_z$ in mm), while $U_x$ and $U_y$ stay within about $\pm 0.5$, so a
relative translation loss almost ignores the 2D position. We instead normalize by the focal length
and the object diameter $D$, as the size-invariant depth of PViT-6D~\cite{stapf2023pvit6d}
(which uses $s = \max(w, h)$):
\begin{equation}
  U_z = \frac{t_z \, s}{\bar f \, D}, \qquad s = \frac{w + h}{2}, \quad
  \bar f = \frac{f_x + f_y}{2}.
  \label{eq:uz}
\end{equation}
Under the pinhole model, an object of diameter $D$ at depth $t_z$ spans about $\bar f D / t_z$
pixels, so $U_z$ is the ratio between the observed box size and the size the object would have if
fully visible: it is dimensionless, close to one, and independent of the camera and of the object
size. At inference, we invert \cref{eq:uz,eq:uxy} with the intrinsics of the test image:
\begin{equation}
  t_z = \frac{U_z \, \bar f \, D}{s}, \quad
  t_x = \frac{(U_x w + b_x - c_x)\, t_z}{f_x}, \quad
  t_y = \frac{(U_y h + b_y - c_y)\, t_z}{f_y}.
  \label{eq:inverse}
\end{equation}
The translation loss is the L2 distance between predicted and target $\U$ relative to the norm of
the target, as in~\cite{ancey2026fastposevit}; with \cref{eq:uz} all three components have the
same scale. We call this formulation \emph{size-invariant} and the original one \emph{legacy}.
\todo{We found this target independently (code name \texttt{site}); credit PViT-6D as prior work
and keep our contribution to the ablation on BOP (\cref{tab:abl_targets}).}

\paragraph{Pixel conventions.}
BOP boxes $[x, y, w, h]$ span the pixels $x$ to $x + w$ inclusive, and pixel $i$ is centred at
$i + 0.5$ in the coordinates of $\K$. We convert every ground-truth and detected box to
$[x, y, x + w + 1, y + h + 1]$ and map it to the crop with one affine warp that accounts for the
half-pixel offsets, so that the box centre lands exactly on the crop centre (checked to
$10^{-5}$ pixels). The crop is the square of side $\max(w, h)$ around the box, so the object keeps
its aspect ratio; areas outside the image are filled with zeros.

\subsection{One network for all objects}
\label{sec:method:conditioning}
A single network must know which object it is looking at, since the same crop content can mean
different poses for different objects (and symmetric objects share appearance). We condition the
shared ViT on the object label $k$ given by the detector with a \emph{per-object class token}: the
[CLS] token of the ViT is replaced by $\mathbf{c} + \mathbf{e}_k$, where $\mathbf{c}$ is the
pretrained class token and $\mathbf{e}_k \in \mathbb{R}^{768}$ a learned embedding of object $k$,
initialized at zero so that training starts from the pretrained network. All objects share the
same transformer layers and the same linear head. Compared with one output head per object
(selecting the nine outputs of object $k$), the token lets the object identity influence every
layer of the transformer. It adds $768$ parameters per object (\cref{tab:size}), and the crops of
different objects of one image form one batch at inference. We compare both options, and one model
per object, in \cref{sec:exp:ablations}.

\subsection{Symmetric objects}
\label{sec:method:symmetry}
For an object with a set of symmetry transformations $\{(\mathbf{S}_j, \mathbf{s}_j)\}$ (from the
BOP model information, continuous symmetries discretized), the poses $(\R\mathbf{S}_j,\ \T +
\R\mathbf{s}_j)$ are all correct. The offset $\mathbf{s}_j$ is zero when the symmetry axis
passes through the model origin, but not always (\eg{} about 19\,mm for the YCB-V clamps). Our
rotation loss takes, per sample, the minimum of the relative Frobenius loss over the symmetric
targets:
\begin{equation}
  \mathcal{L}_R = \min_j \frac{\| \mathbf{r}_{6D} - \mathrm{6D}(\R'_j) \|}
  {\| \mathrm{6D}(\R'_j) \|},
  \label{eq:symloss}
\end{equation}
where $\R'_j$ is the apparent rotation of the $j$-th symmetric pose. Since the symmetries act on
the right (model frame) and the apparent-rotation correction on the left (camera frame), the
symmetric targets are obtained directly from the apparent rotation when $\mathbf{s}_j = 0$. When
$\mathbf{s}_j \neq 0$, the symmetric pose has its own translation, hence its own $\U$ and its own
correction $\Delta\R$: we build the labels of every symmetric pose from the original ground truth,
take the minimum of \cref{eq:symloss}, and use the translation label of the chosen symmetry in the
translation loss.

\subsection{Training data and augmentations}
\label{sec:method:training}

\paragraph{Data.}
We train only on the PBR images provided by BOP~\cite{denninger2019blenderproc}: for each
dataset, the instances of the test objects with a visible fraction of at least 10\,\%, cropped
around their visible bounding box (\texttt{bbox\_visib}), which is what a detector sees. Two PBR
scenes are held out for validation. Objects are sampled in balanced epochs: every epoch draws the
same number of instances per object, so frequent objects do not dominate.

\paragraph{Augmentations.}
Pixel-level augmentations are applied independently with probability 0.5 each, with
Albumentations: brightness and contrast, gamma, hue and saturation, gray, tone curve, CLAHE,
solarize, sun flare, Gaussian, motion and defocus blur, Gaussian and ISO noise, downscaling and
JPEG compression (the ``domain gap'' set of~\cite{ancey2026fastposevit} without its geometric
distortion, which would change the pose labels). The bounding box is perturbed with probability
0.5 by up to 10\,\% of its size on each side. We do not use the in-plane rotation augmentation of
the original method: rotating the image about its centre is a pure camera roll only if the
principal point is at the image centre and $f_x = f_y$, which does not hold for the BOP cameras
(on LM, it would cause label errors of about 20\,mm at 180\textdegree).
\pending{Augmentation study (\cref{sec:exp:augmentations}) may change this paragraph.}

\subsection{Multi-class detector}
\label{sec:method:detector}
For a fully real-time pipeline, we train one LW-DETR~\cite{chen2024lwdetr} for all objects of a
dataset on the same PBR images, with the original LW-DETR recipe (group DETR with 13 query
groups, IoU-aware classification loss, EMA, AdamW, COCO-pretrained weights except the class head)
and the same pixel augmentations plus horizontal flips, at $640 \times 640$ pixels. We train for
50 epochs over the 48,000 PBR images of LM (batch size 32, learning rate $10^{-4}$, $1.5 \cdot
10^{-4}$ for the encoder, weight decay $10^{-4}$, 500 warm-up iterations, EMA decay 0.9997,
gradient clipping at 0.1) and keep the checkpoint with the best COCO mAP on the held-out PBR
scenes. Training takes about 12\,h (``tiny'' variant) and 22\,h (``small'') on one A100.
